\begin{papercard}{apa2024}
\cf{Stages} \stagetag{A} (Step 1 background attributes and Step 3 privacy questions as sampling and screening attributes), \stagetag{C} (question list as seed for codebook categories), \stagetag{H} (the framework's list of digital risks), \stagetag{S} (privacy and security placed as a gate before clinical evaluation). A 6-page print (dated 20 October 2026, PDF footer ``printed 10/20/26'') of the American Psychiatric Association's web page ``The App Evaluation Model'' (psychiatry.org, App Advisor). It is a clinician-facing, unscored framework for choosing mental health apps with a patient, in six hierarchical steps. It is a web page, not a peer-reviewed paper: no authors, no version number and no validation are given in the file.

\cf{Stance and purpose} The page rests on the belief that choosing an app is ``a personal decision based on many factors'' with rarely a binary yes or no, and notes that picking an app is not something psychiatrists are classically trained to do \pc{apa2024}{1}{Our approach to rating mental health apps is grounded in the belief that any decision between you and a patient is a personal decision}. The goal is ``a hierarchical rating system with content-specific questions that guide a thoughtful approach to app selection'' \pc{apa2024}{1}{provide a hierarchical rating system with content-specific questions that guide a thoughtful approach to app selection}. There is no minimum or maximum number of levels an app must meet to count as ``good'' or ``useful''; the model shows presence and absence of features to support an informed choice \pc{apa2024}{2}{There is no minimum or maximum number of levels of the model the app needs to meet to be considered}.

\cf{The six steps} (1) \emph{Background}: ownership, trusted source, funding and conflicts of interest, stated purpose, medical claims, costs, ``updated in the last 180 days'', feedback channel \pc{apa2024}{2}{Has the app been updated in the last 180 days}; one need not answer every question to proceed. (2) \emph{Access}: low-bandwidth or offline use, platforms, web use, accessibility features, languages, reading level, support needs, intended users \pc{apa2024}{3}{Does the app function in low bandwidth or offline environments}. (3) \emph{Privacy and security} (below). (4) \emph{Clinical foundation}: accuracy, clinical or recovery basis, evidence sources, studies among intended users, retention rates \pc{apa2024}{5}{Does the app disclose clinical retention rates}, user involvement in design. (5) \emph{Usability}: modes of engagement, fit with priorities, customisation, navigation, support, personalisation. (6) \emph{Integration toward patient-centred goals}: individual versus team use, data access and export \pc{apa2024}{6}{Can the user export/download their data for their own use}, secure sharing with care teams, electronic medical records and wearables.

\cf{Step 3: privacy and security} The text says ``digital tools present unique risks that are often underestimated'' and lists exposure to ``data costs, social profiling, loss of insurance benefits, or breaches of sensitive health information'' \pc{apa2024}{3}{Apps may expose users to data costs, social profiling, loss of insurance benefits, or breaches of sensitive health information}. The questions are ``not exhaustive, and there is currently no gold standard for rating these features''; most answers should be in the privacy policy or on the developer's site, and ``if no clear policy exists, the app may not be appropriate for handling sensitive health information'' \pc{apa2024}{3}{If no clear policy exists, the app may not be appropriate for handling sensitive health information}. Some details can only be taken from the developer and ``you may need to take these at face value''; there is no scoring system for the level \pc{apa2024}{4}{Because some details, such as security measures, may only be described by the developer, you may need to take these at face value}. The 11 questions: easy-to-understand policy; reading level; declared data use and purpose; collection, use or transmission of sensitive data such as protected health information, and de-identification; secure authentication; opt-out and deletion, including whether the process is stated; storage location; third parties receiving data \pc{apa2024}{4}{What third parties does the app share data with}; response to clinical safety concerns; a plan for data breaches and how past incidents were handled \pc{apa2024}{4}{Does the developer disclose a plan for addressing breaches of data security}; and whether AI is used and how \pc{apa2024}{4}{Does the app use AI? If so, does it describe how AI is used}, with monitoring and mitigation of bias.

\cf{Gating logic} Step 4 begins ``if you decide that an app has sufficient privacy and security'' \pc{apa2024}{4}{App developers often make many claims on their apps' clinical effectiveness or background}; the clinician is advised to download and try the app \pc{apa2024}{4}{we recommend that you download and try the app to see what it is actually doing}. By Step 5 an app that passed Steps 1--4 ``offers minimal risk in terms of digital safety and privacy'' \pc{apa2024}{5}{It offers minimal risk in terms of digital safety and privacy}, and data integration only becomes meaningful after the app is usable, safe, evidence-informed and user-friendly \pc{apa2024}{5}{Data sharing only becomes meaningful once the app has been deemed usable}. Usability is called ``a more subjective category'' \pc{apa2024}{5}{usability is a more subjective category}.

\cf{Limits (card)} A professional-body web page: no method for developing the questions, no reliability or validity data, no scoring, no examples of rated apps in the file; answers depend on developer self-report; the claim that an app passing Steps 1--4 has ``minimal risk'' is an assumption not backed by any check of actual data flows; US psychiatry audience; the page layout lists the six levels as a diagram whose first label is misspelled; the content may change after the print date, so it should be cited with an access date. The framework treats privacy questions as yes/no presence of statements, so it cannot show how harms are \emph{communicated}.

\cf{Use for this thesis} (1) Evidence that clinician guidance places privacy and security before benefits and asks about the same things the thesis codes in policies: purpose, sensitive data, de-identification, opt-out and deletion, storage location, third parties, breach plan, AI. The 11 questions are a cross-check for codebook completeness (stage C). (2) Step 1 and Step 2 attributes (ownership, funding, update within 180 days, cost, platforms, languages, reading level) are ready-made descriptive fields for the app sheet (stages A and D); ``updated in the last 180 days'' is a concrete recency criterion that can be reported and justified. (3) The list of risks (data costs, social profiling, loss of insurance benefits, breaches) can seed harm-tree branches, but it is unreferenced on the page and should be labelled as a practitioner statement. (4) The admission that no gold standard exists and that developer statements must be taken at face value supports the thesis' stance that it measures what apps \emph{say}, not what they do. (5) Cite with the access date 20 October 2026 and note that it is a web page; pair with Torous and Henson papers if added from the candidates list.
\end{papercard}
